% !TeX program = lualatex
% !TeX encoding = utf8
% !TeX spellcheck = uk_UA
% !TEX indexprogram = upmendex
% !TeX root =./main.tex

%=========================================================
\chapter*{Передмова}\label{\currfilebase}\hypertarget{\currfilebase}{}
\addcontentsline{toc}{chapter}{Передмова}
%=========================================================


Переважна більшість баріонної речовини у~Всесвіті~--- зорі, міжзоряний газ, сонячний
вітер~--- перебуває в~стані плазми. На Землі плазма~--- це блискавка, полярні сяйва,
іоносфера, газовий розряд у~лампах, зварювальна дуга, плазмове травлення мікросхем
і~плазма, яку намагаються утримувати в~установках керованого термоядерного синтезу.

Спільна риса всіх цих явищ~--- колективна поведінка заряджених частинок. У~всьому розділі
важливо відрізняти \emphx{масштабну оцінку} від точної теорії. Якщо результат отримано
в~простій моделі, ми спочатку з'ясуємо фізичний зміст цієї моделі, а~потім зазначимо, де
вона перестає працювати.




Фізика плазми поєднує кілька рівнів опису. Спочатку ми
визначимо, коли система заряджених частинок справді поводиться як плазма, і~введемо її
характерні просторові та часові масштаби. Потім розглянемо рух \emph{окремої} частинки
в~електромагнітному полі, зіткнення та перенос. Після цього перейдемо до
\emph{колективних} коливань і~хвиль, а~далі --- до гідродинамічного та
магнітогідродинамічного опису. Нарешті застосуємо ці моделі до магнітного утримання
і~керованого термоядерного синтезу.


\emphx{Головний принцип посібника:} одна й~та сама плазма може потребувати різних моделей
залежно від масштабу явища. Тому щоразу будемо ставити три питання: \emph{який просторовий
масштаб? який часовий масштаб? які фізичні процеси встигають відбутися?}


%% --------------------------------------------------------
\section*{Система одиниць}
%% --------------------------------------------------------

У~цьому посібнику, як і~в~більшості наведених нижче формул, використано систему СГС-Гаусса:
сила Лоренца має вигляд $\vect F=q(\vect E+\vect v\times\vect B/c)$, а~закон Ампера ---
$\rot\vect B=(4\pi/c)\vect j+(1/c)\partial_t\vect E$. Наприкінці розділу наведено
практичні формули в~СГС. Перехід до СІ простий, але множники $c$ та $4\pi$ в~окремих
рівняннях змінюються.